problem:  
Let \((X,d,\mathfrak{m})\) be a **noncollapsed** \(\mathrm{RCD}(0,N)\) space with \(1<N<\infty\), i.e., \(\mathfrak{m}=\mathcal{H}^N\).  
For \(p\in(1,N)\), define the sharp Sobolev constant
\[
C_p(X):=\sup_{f\in W^{1,p}(X)\setminus\{0\}}
\frac{\|f\|_{L^{p^*}(\mathfrak{m})}^{p}}{\|Df\|_{L^{p}(\mathfrak{m})}^{p}}, 
\quad p^*=\frac{Np}{N-p}.
\]
On Euclidean space \(\mathbb{R}^N\), denote the classical constant by \(C_p^{\mathrm{Euc}}\), and set
\[
A_p^{\mathrm{Euc}}:=(N-p)^{p/N}C_p^{\mathrm{Euc}},
\]
which converges to a finite limit determining the critical Moser–Trudinger constant.  

**Open Problem (Critical Sobolev asymptotics and Euclidean rigidity):**  

(a) *(Existence of limit)* For general noncollapsed \(\mathrm{RCD}(0,N)\) spaces, does the renormalized sequence  
\[
A_p(X):=(N-p)^{p/N}C_p(X)
\]
admit a finite limit as \(p\to N^{-}\)?  

(b) *(Rigidity at the Euclidean value)* If the limit exists, does
\[
\lim_{p\to N^-}\frac{A_p(X)}{A_p^{\mathrm{Euc}}}=1
\]
hold **if and only if** \(X\) is globally isometric (up to measure scaling) to the Euclidean space \((\mathbb{R}^N,|\cdot|,c\,dx)\)?  

(c) *(Dependence on local geometry)*  
For \(\mathcal{H}^N\)-a.e. regular point \(x\in X\), the infinitesimal volume density
\[
\theta(x):=\lim_{r\to0^+}\frac{\mathfrak{m}(B_r(x))}{\omega_N r^N}
\]
exists and may vary with \(x\).  
Does the asymptotic ratio
\[
L(X):=\lim_{p\to N^-}\frac{A_p(X)}{A_p^{\mathrm{Euc}}}
\]
depend only on the **distribution** of \(\theta(x)\) (for instance, on its essential infimum or mean value), thereby defining a geometric analytic invariant linking Sobolev asymptotics to infinitesimal volume structure?

Why is it a "good" problem:  
This formulation corrects earlier issues and is now mathematically sound. It focuses on the **limit behavior of sharp Sobolev constants** at the critical exponent, within the synthetic curvature framework, and connects it with the local measure–metric geometry encoded by volume densities.  

This makes it a *good* problem because:  
- It isolates a clean analytic phenomenon—the asymptotics of \(C_p(X)\) as \(p\to N^-\)—and proposes a deep geometric interpretation.  
- It builds a bridge between **nonlinear analysis** (sharp critical inequalities and \(p\)-Laplacian asymptotics) and **synthetic geometry** (tangent cones, measure densities, and noncollapsing structure).  
- The conjectured rigidity at the Euclidean value parallels Obata-type theorems but in a limit and nonsmooth regime, potentially introducing new analytic invariants for metric spaces.  

The problem is simple to state yet probes profound interactions between curvature, analysis, and geometry, embodying the essence of a forward-looking research-level question.